\documentclass[12pt]{article}
\usepackage{graphicx} % Required for inserting images
\usepackage[margin=0.5in]{geometry}
\usepackage{amsmath, amssymb}
\usepackage{mathrsfs}


\usepackage{soul}


% Dr. David Brown Commands
\newcommand{\ds}{\displaystyle}
\newcommand{\R}{\mathbb{R}}
\newcommand{\N}{\mathbb{N}}
\newcommand{\Q}{\mathbb{Q}}
\newcommand{\C}{\mathbb{C}}


% Commands to get script r
\usepackage{calligra}
\DeclareMathAlphabet{\mathcalligra}{T1}{calligra}{m}{n}
\DeclareFontShape{T1}{calligra}{m}{n}{<->s*[2.2]callig15}{}
\newcommand{\scripty}[1]{\ensuremath{\mathcalligra{#1}}}

% Unit commands
\newcommand{\degree}{\text{°}}
\newcommand{\unitSpeed}{\frac{\text{m}}{\text{s}}}
\newcommand{\unitAcceleration}{\frac{\text{m}}{\text{s}^2}}

% Constant commands
\newcommand{\avogadro}{6.022\times10^{23}}

\newcommand{\boltzmannConstK}{1.381\times10^{-23}\frac{\text{J}}{\text{K}}}
\newcommand{\planckConst}{6.626\times10^{-34}\text{J}\cdot\text{s}}

\newcommand{\atmP}{1.01325\times10^5\,\text{Pa}}

% Commands to easily write partial derivatives
\newcommand{\partDeriv}[2]{\frac{\partial #1}{\partial #2}}
\newcommand{\doublePartDeriv}[2]{\frac{\partial^2 #1}{\partial^2 #2}}


\title{Problem 1.20}
\author{Gabriel Decker}


\begin{document}



\maketitle
\begin{flushleft}

For this problem, we need the following formula.
\begin{equation}
    v_\text{rms}=\sqrt{\frac{3kT}{m}}
\end{equation}
This is the root-mean-squared speed of an ideal gas.
The atomic mass of Fluorine is 18.998 amu.
By the conversion factor of $1\,\text{amu}=1.6605\times10^{-27}\,\text{kg}$, this implies that the mass of a Fluroine atom is $m_F=3.155\times10^{-26}\,\text{kg}$.
This makes the mass of six Fluorine atoms (some mass will be lost when they bond, but assume that's negligible) $m_{F6}=1.893\times10^{-25}\,\text{kg}$.
The atomic weight of Uranium-238 is about 238 amu or $m_{U238}=3.952\times10^{-25}\,\text{kg}$.
The atomic weight of Uranium-235 is about 235 amu or $m_{U238}=3.902\times10^{-25}\,\text{kg}$.\\~\\

Therefore, the atomic weights of the two molecules are about $m_1=5.8449\times10^{-25}\,\text{kg}$ for Uranium-238 and $m_2=5.7951\times10^{-25}\,\text{kg}$ for Uranium-235.
Since room temp is about $T=300\,\text{K}$, this implies that
$$v_\text{1 rms}=\sqrt{\frac{3\cdot(\boltzmannConstK)\cdot(300\,\text{K})}{5.8449\times10^{-25}\,\text{kg}}}$$
$$\implies v_\text{1 rms}=145.82\,\unitSpeed$$
and
$$v_\text{2 rms}=\sqrt{\frac{3\cdot(\boltzmannConstK)\cdot(300\,\text{K})}{5.7951\times10^{-25}\,\text{kg}}}$$
$$\implies v_\text{2 rms}=146.45\,\unitSpeed$$




\end{flushleft}

\end{document}
